% Options for packages loaded elsewhere
\PassOptionsToPackage{unicode}{hyperref}
\PassOptionsToPackage{hyphens}{url}
\PassOptionsToPackage{dvipsnames,svgnames,x11names}{xcolor}
\documentclass[
  a4paper,
]{article}
\usepackage{xcolor}
\usepackage[margin=20mm]{geometry}
\usepackage{amsmath,amssymb}
\setcounter{secnumdepth}{-\maxdimen} % remove section numbering
\usepackage{iftex}
\ifPDFTeX
  \usepackage[T1]{fontenc}
  \usepackage[utf8]{inputenc}
  \usepackage{textcomp} % provide euro and other symbols
\else % if luatex or xetex
  \usepackage{unicode-math} % this also loads fontspec
  \defaultfontfeatures{Scale=MatchLowercase}
  \defaultfontfeatures[\rmfamily]{Ligatures=TeX,Scale=1}
\fi
\usepackage{lmodern}
\ifPDFTeX\else
  % xetex/luatex font selection
  \setmainfont[]{DejaVu Serif}
  \setsansfont[]{DejaVu Sans}
  \setmonofont[]{DejaVu Sans Mono}
\fi
% Use upquote if available, for straight quotes in verbatim environments
\IfFileExists{upquote.sty}{\usepackage{upquote}}{}
\IfFileExists{microtype.sty}{% use microtype if available
  \usepackage[]{microtype}
  \UseMicrotypeSet[protrusion]{basicmath} % disable protrusion for tt fonts
}{}
\makeatletter
\@ifundefined{KOMAClassName}{% if non-KOMA class
  \IfFileExists{parskip.sty}{%
    \usepackage{parskip}
  }{% else
    \setlength{\parindent}{0pt}
    \setlength{\parskip}{6pt plus 2pt minus 1pt}}
}{% if KOMA class
  \KOMAoptions{parskip=half}}
\makeatother
\usepackage{longtable,booktabs,array}
\newcounter{none} % for unnumbered tables
\usepackage{calc} % for calculating minipage widths
% Correct order of tables after \paragraph or \subparagraph
\usepackage{etoolbox}
\makeatletter
\patchcmd\longtable{\par}{\if@noskipsec\mbox{}\fi\par}{}{}
\makeatother
% Allow footnotes in longtable head/foot
\IfFileExists{footnotehyper.sty}{\usepackage{footnotehyper}}{\usepackage{footnote}}
\makesavenoteenv{longtable}
\usepackage{graphicx}
\makeatletter
\newsavebox\pandoc@box
\newcommand*\pandocbounded[1]{% scales image to fit in text height/width
  \sbox\pandoc@box{#1}%
  \Gscale@div\@tempa{\textheight}{\dimexpr\ht\pandoc@box+\dp\pandoc@box\relax}%
  \Gscale@div\@tempb{\linewidth}{\wd\pandoc@box}%
  \ifdim\@tempb\p@<\@tempa\p@\let\@tempa\@tempb\fi% select the smaller of both
  \ifdim\@tempa\p@<\p@\scalebox{\@tempa}{\usebox\pandoc@box}%
  \else\usebox{\pandoc@box}%
  \fi%
}
% Set default figure placement to htbp
\def\fps@figure{htbp}
\makeatother
\setlength{\emergencystretch}{3em} % prevent overfull lines
\providecommand{\tightlist}{%
  \setlength{\itemsep}{0pt}\setlength{\parskip}{0pt}}
\usepackage{luatexja}
\usepackage{xurl}
\usepackage{float}
\usepackage{bookmark}
\IfFileExists{xurl.sty}{\usepackage{xurl}
\usepackage{float}}{} % add URL line breaks if available
\urlstyle{same}
\hypersetup{
  pdftitle={Spherically Symmetric Spatial Metric Distortion{,} Discrete Curvature Gradient{,} Spin-2 Clock Dragging{,} and Critical Mass Collapse in the 3D Triangle-Loop Product (Fifth Rapid Communication)},
  pdfauthor={Noriaki Kihara (WF System Co., Ltd.) ORCID 0009-0004-6753-4020},
  colorlinks=true,
  linkcolor={Maroon},
  filecolor={Maroon},
  citecolor={Blue},
  urlcolor={Blue},
  pdfcreator={LaTeX via pandoc}}

\title{Spherically Symmetric Spatial Metric Distortion, Discrete
Curvature Gradient, Spin-2 Clock Dragging, and Critical Mass Collapse in
the 3D Triangle-Loop Product (Fifth Rapid Communication)}
\usepackage{etoolbox}
\makeatletter
\providecommand{\subtitle}[1]{% add subtitle to \maketitle
  \apptocmd{\@title}{\par {\large #1 \par}}{}{}
}
\makeatother
\subtitle{― Exact 7-Direction Isotropic 4-Class Metric Splitting,
Far-Space Disconnection at \(\epsilon_c \in (0.270, 0.275)\), and
Rigorous Phase-Rigidity Audit ―}
\author{Noriaki Kihara (WF System Co., Ltd.) ORCID 0009-0004-6753-4020}
\date{November 28, 2026}

\begin{document}
\maketitle

\textbf{Author:} Noriaki Kihara (WF System Co., Ltd.)\\
\textbf{ORCID:} 0009-0004-6753-4020\\
\textbf{Date:} November 28, 2026\\
\textbf{Status:} Fifth Rapid Communication following Rapid Reports 1--4
(DOIs: 10.5281/zenodo.24616931, 24618400, 24619357, 24620262).
Quantifies the spatial metric distortion, discrete curvature gradient,
critical mass collapse, and spin-2 clock dragging on the Layer-1 time
crystal \(H = O_8 \otimes M_8\), while explicitly reporting a rigorous
falsification audit showing single-clock phase rigidity in the
weak-coupling background.

\begin{center}\rule{0.5\linewidth}{0.5pt}\end{center}

\subsection{Abstract}\label{abstract}

We investigate the geometric and temporal response of the Layer-1 time
crystal \(H = O_8 \otimes M_8\) in the background-free \(N = 64\) real
3-tensor triangle-loop product
\(\mathrm{Tri}(T)_{ijk} = \sum_{p,q,r} T_{ipq} T_{qjr} T_{rpk}\) when
perturbed by a single-site localized condensate P4 and a spin-2 mode
(\(\sigma_x\)), using deterministic zero-randomness simulations:

\begin{enumerate}
\def\labelenumi{\arabic{enumi}.}
\tightlist
\item
  \textbf{Spherically Symmetric 4-Class Spatial Metric Distortion
  (7-Digit Isotropic)}: Placing P4 (\(\epsilon = 0.16, \tau = 40\)) at
  site \(A = 0\) splits the 64 background \(O_8\) edges into four exact
  geometric classes with 7-digit isotropy across all 7 imaginary
  directions \(m = 1 \dots 7\): center self-loop \((0,0,0)\) at
  \(E/E_{\mathrm{ref}} = 1.3248375\) (\(+32.5\%\)), far tangential edges
  \((m,n,p)\) at \(1.0718855\) (\(+7.2\%\)), radial connecting edges
  \((0,m,m), (m,m,0)\) at \(1.0122646\) (\(+1.2\%\)), and far self-loops
  \((m,0,m)\) at \(1.0000000\). Control tests at imaginary site
  \(A = 1\) and symmetric two-mass configurations (\(A = 0, 1\): 6-digit
  match; remaining 6 sites uniform to \(7.5 \times 10^{-17}\)) confirm
  full site equivalence.
\item
  \textbf{Discrete Curvature (\(\rho = L_3/L_1\)) Gradient}: While
  vacuum \(H\) is strictly flat across all sites
  (\texttt{n\_tri\ =\ 64}, weighted view-angle
  \(\langle\rho\rangle_w = 0.576699\)), inserting P4
  (\(\epsilon = 0.20, \tau = 30\)) adds \(+8\) cross triangles
  (\texttt{n\_tri\ =\ 72}) and generates a spherically symmetric
  curvature gradient with \(\langle\rho\rangle_w = 0.796889\) at mass
  site \(A = 0\) versus \(0.787796\) at far sites \(A = 1, 7\).
\item
  \textbf{Critical Mass Collapse at \(\epsilon_c \in (0.270, 0.275)\)}:
  Crossing \(\epsilon_c\) triggers a sharp first-order collapse
  (\texttt{nz:\ 576\ \$\textbackslash{}to\$\ 16},
  \(\|P\|: 0.04682 \to 0.49998 = 1/2\)) in which the radial links sever
  and all far-space energies
  \(E_{\mathrm{bg}}(\mathrm{far}), E_P(\mathrm{far})\) vanish to
  \texttt{0.000000}, while background (\texttt{0.2768}) and particle
  (\texttt{0.7232}) sectors remain coexisting inside the single mass
  site. Collapsed end-states at \(A = 0\) and \(A = 1\)
  (\(\epsilon = 0.30\)) match to 6 digits.
\item
  \textbf{Rigorous Phase-Rigidity Audit vs.~Spin-2 Clock Dragging}: A
  direct complex-phase audit reveals that all background edge phases
  remain synchronized to \(|\Delta\theta| < 10^{-9}\) rad across
  \(\tau = 1 \dots 100\) (ruling out cumulative background frequency
  drift in weak coupling and identifying an apparent time-averaged ratio
  \(\delta\phi/\delta g \approx 0.71\) as a reference-angle artifact).
  Conversely, genuine temporal slowdown occurs in the particle sector
  (\(\phi_P/\phi_0 = 0.56\)--\(0.86\)) and under spin-2 (\(\sigma_x\))
  coupling, which universally locks to \(d_1 = 1/\sqrt{2}\)
  (\(\phi/\phi_0 = 0.6578\)) and drags the entire background clock down
  to \texttt{0.6579} at \(\epsilon = 0.20\).
\end{enumerate}

\begin{center}\rule{0.5\linewidth}{0.5pt}\end{center}

\subsection{1. Introduction and Audit
Methodology}\label{introduction-and-audit-methodology}

Rapid Communications 1--4 {[}1--4{]} established Hurwitz crystal
selection, four-stage doubling-entanglement cascades, \(SO(3)\) spin
selection, \(G_2 \supset SU(3)\) gauge counts, \(\mathbb{Z}_2^3\)
spatial harmonic localization, two-particle coexistence, and four
emergent quantum theorems (Q1--Q4). Here we probe how localized mass and
spin-2 excitations deform the spatial geometry and clock rates of the
Layer-1 background \(H = O_8 \otimes M_8\). Because gravitational claims
are prone to measurement artifacts, we pair main sweeps
(\texttt{tri\_loop\_n64\_gravity\_time\_lorentz\_v0.py}) with a
five-point falsification audit
(\texttt{tri\_loop\_n64\_gravity\_rigorous\_check\_v0.py}) and complete
edge-by-edge dissection
(\texttt{tri\_loop\_n64\_dphi\_dg\_anatomy\_v0.py}).

\begin{center}\rule{0.5\linewidth}{0.5pt}\end{center}

\subsection{2. Spherically Symmetric 4-Class Spatial Metric
Distortion}\label{spherically-symmetric-4-class-spatial-metric-distortion}

Adding P4 (\(\delta_{000} \otimes M_8[b=\{2,3\}]\)) at \(A = 0\) and
evolving to \(\tau = 40\) splits all 64 background (\(b = 0\)) \(O_8\)
edges \((A, B, C)\) into four exact geometric classes normalized by
\(E_{\mathrm{ref}} = E(1, 0, 1)\) (Table 1, Fig. 1(a)):

{\def\LTcaptype{none} % do not increment counter
\begin{longtable}[]{@{}
  >{\raggedright\arraybackslash}p{(\linewidth - 8\tabcolsep) * \real{0.1739}}
  >{\raggedright\arraybackslash}p{(\linewidth - 8\tabcolsep) * \real{0.1739}}
  >{\centering\arraybackslash}p{(\linewidth - 8\tabcolsep) * \real{0.2174}}
  >{\centering\arraybackslash}p{(\linewidth - 8\tabcolsep) * \real{0.2174}}
  >{\centering\arraybackslash}p{(\linewidth - 8\tabcolsep) * \real{0.2174}}@{}}
\toprule\noalign{}
\begin{minipage}[b]{\linewidth}\raggedright
Spatial Edge Class
\end{minipage} & \begin{minipage}[b]{\linewidth}\raggedright
Components (Count)
\end{minipage} & \begin{minipage}[b]{\linewidth}\centering
\(E/E_{\mathrm{ref}}\) (\(\epsilon=0.08\))
\end{minipage} & \begin{minipage}[b]{\linewidth}\centering
\(E/E_{\mathrm{ref}}\) (\(\epsilon=0.16\))
\end{minipage} & \begin{minipage}[b]{\linewidth}\centering
Phase Diff \(\|\Delta\theta\|\)
\end{minipage} \\
\midrule\noalign{}
\endhead
\bottomrule\noalign{}
\endlastfoot
\textbf{1. Center self-loop} & \((0, 0, 0)\) (1) & \texttt{1.1173232}
(\(+11.7\%\)) & \textbf{\texttt{1.3248375} (\(+32.5\%\))} &
\(6.07 \times 10^{-11}\) rad \\
\textbf{2. Radial links} & \((0, m, m), (m, m, 0)\) (14) &
\texttt{1.0054810} (\(+0.55\%\)) & \textbf{\texttt{1.0122646}
(\(+1.23\%\))} & \(6.08 \times 10^{-9}\) rad \\
\textbf{3. Far tangential} & \((m, n, p)\; (m,n,p\neq 0)\) (42) &
\texttt{1.0257321} (\(+2.57\%\)) & \textbf{\texttt{1.0718855}
(\(+7.19\%\))} & \(6.10 \times 10^{-9}\) rad \\
\textbf{4. Far self-loops} & \((m, 0, m)\) (7) & \texttt{1.0000000}
(ref) & \texttt{1.0000000} (ref) & \(0\) (ref) \\
\end{longtable}
}

All 7 imaginary directions \(m = 1 \dots 7\) are strictly isotropic to 7
digits, with zero leakage into \(O_8 = 0\) entries
(\texttt{0.000000e+00}). Energy concentrates strongly at the center
\((0,0,0)\) while radial connecting links \((0,m,m)\) are relatively
suppressed compared to far tangential edges. Control experiments placing
P4 at imaginary site \(A = 1\) reproduce the same distortion
(\texttt{dg\ =\ 0.002967} at \(\epsilon = 0.16\)), and two-mass
placement (\(A = 0, 1\), \(\epsilon = 0.12\) each) yields a 6-digit
symmetric well (\texttt{dg\ =\ 0.003451} at both sites) with the
remaining 6 sites uniform to \(7.52 \times 10^{-17}\).

\begin{center}\rule{0.5\linewidth}{0.5pt}\end{center}

\subsection{\texorpdfstring{3. Discrete Curvature (\(\rho = L_3/L_1\))
Gradient}{3. Discrete Curvature (\textbackslash rho = L\_3/L\_1) Gradient}}\label{discrete-curvature-rho-l_3l_1-gradient}

Measuring the triangle-loop view-angle \(\rho = L_3/L_1\) {[}1, 2{]} at
mass site \(A = 0\) versus far sites \(A = 1, 7\) confirms a transition
from exact flatness in vacuum \(H\) to a spherically symmetric discrete
curvature gradient (Table 2, Fig. 1(b)):

{\def\LTcaptype{none} % do not increment counter
\begin{longtable}[]{@{}
  >{\raggedright\arraybackslash}p{(\linewidth - 8\tabcolsep) * \real{0.1667}}
  >{\centering\arraybackslash}p{(\linewidth - 8\tabcolsep) * \real{0.2083}}
  >{\centering\arraybackslash}p{(\linewidth - 8\tabcolsep) * \real{0.2083}}
  >{\centering\arraybackslash}p{(\linewidth - 8\tabcolsep) * \real{0.2083}}
  >{\centering\arraybackslash}p{(\linewidth - 8\tabcolsep) * \real{0.2083}}@{}}
\toprule\noalign{}
\begin{minipage}[b]{\linewidth}\raggedright
Configuration / Site
\end{minipage} & \begin{minipage}[b]{\linewidth}\centering
Triangles \(n_{\mathrm{tri}}\)
\end{minipage} & \begin{minipage}[b]{\linewidth}\centering
Mean \(\bar{\rho}\)
\end{minipage} & \begin{minipage}[b]{\linewidth}\centering
Weighted \(\langle\rho\rangle_w\)
\end{minipage} & \begin{minipage}[b]{\linewidth}\centering
Min \(\rho_{\min}\)
\end{minipage} \\
\midrule\noalign{}
\endhead
\bottomrule\noalign{}
\endlastfoot
Vacuum \(H\) (\(A = 0, 1, 7\)) & 64 & \texttt{0.621227} &
\texttt{0.576699} & \texttt{0.414735} \\
With Mass (\(\epsilon=0.20\)): Far \(A = 1, 7\) & 72 (\(+8\)) &
\texttt{0.709985} & \texttt{0.787796} & \texttt{0.064441} \\
\textbf{With Mass (\(\epsilon=0.20\)): Center \(A = 0\)} & \textbf{72
(\(+8\))} & \textbf{\texttt{0.716065}} & \textbf{\texttt{0.796889}} &
\texttt{0.064441} \\
\end{longtable}
}

\begin{figure}[H]
\centering
\includegraphics[width=0.96\linewidth,height=\textheight,keepaspectratio,alt={(a) Spherically symmetric 4-class spatial metric distortion induced by P4 (\textbackslash epsilon=0.08, 0.16, 7-digit isotropic). (b) Discrete curvature (\textbackslash rho = L\_3/L\_1) gradient comparing flat vacuum H with center A=0 and far sites A=1,7.}]{figs_en/fig_r5_metric_and_curvature_en.png}
\caption{(a) Spherically symmetric 4-class spatial metric distortion
induced by P4 (\(\epsilon=0.08, 0.16\), 7-digit isotropic). (b) Discrete
curvature (\(\rho = L_3/L_1\)) gradient comparing flat vacuum \(H\) with
center \(A=0\) and far sites \(A=1,7\).}
\end{figure}

\begin{center}\rule{0.5\linewidth}{0.5pt}\end{center}

\subsection{\texorpdfstring{4. Critical Mass Collapse at
\(\epsilon_c \in (0.270, 0.275)\)}{4. Critical Mass Collapse at \textbackslash epsilon\_c \textbackslash in (0.270, 0.275)}}\label{critical-mass-collapse-at-epsilon_c-in-0.270-0.275}

Scanning \(\epsilon \in [0.250, 0.300]\) in steps of \texttt{0.005}
(\(\tau = 60\)) pinpoints a sharp first-order collapse at
\(\epsilon_c \in (0.270, 0.275)\) (Table 3, Fig. 2(a)):

{\def\LTcaptype{none} % do not increment counter
\begin{longtable}[]{@{}
  >{\centering\arraybackslash}p{(\linewidth - 14\tabcolsep) * \real{0.1250}}
  >{\centering\arraybackslash}p{(\linewidth - 14\tabcolsep) * \real{0.1250}}
  >{\centering\arraybackslash}p{(\linewidth - 14\tabcolsep) * \real{0.1250}}
  >{\centering\arraybackslash}p{(\linewidth - 14\tabcolsep) * \real{0.1250}}
  >{\centering\arraybackslash}p{(\linewidth - 14\tabcolsep) * \real{0.1250}}
  >{\centering\arraybackslash}p{(\linewidth - 14\tabcolsep) * \real{0.1250}}
  >{\centering\arraybackslash}p{(\linewidth - 14\tabcolsep) * \real{0.1250}}
  >{\centering\arraybackslash}p{(\linewidth - 14\tabcolsep) * \real{0.1250}}@{}}
\toprule\noalign{}
\begin{minipage}[b]{\linewidth}\centering
\(\epsilon\)
\end{minipage} & \begin{minipage}[b]{\linewidth}\centering
\texttt{nz}
\end{minipage} & \begin{minipage}[b]{\linewidth}\centering
\(\|P\|\)
\end{minipage} & \begin{minipage}[b]{\linewidth}\centering
\(E_{\mathrm{bg}}(0)\)
\end{minipage} & \begin{minipage}[b]{\linewidth}\centering
\(E_{\mathrm{bg}}(\mathrm{far})\)
\end{minipage} & \begin{minipage}[b]{\linewidth}\centering
\(E_P(0)\)
\end{minipage} & \begin{minipage}[b]{\linewidth}\centering
\(E_P(\mathrm{far})\)
\end{minipage} & \begin{minipage}[b]{\linewidth}\centering
SVD \(\sigma_1, \sigma_2\)
\end{minipage} \\
\midrule\noalign{}
\endhead
\bottomrule\noalign{}
\endlastfoot
\texttt{0.260} & 576 & \texttt{0.04682} & \texttt{0.124530} &
\texttt{0.875458} & \texttt{0.000002} & \texttt{0.000010} &
\texttt{1.0000,\ 0.0032} \\
\texttt{0.270} & 576 & \texttt{0.04682} & \texttt{0.124188} &
\texttt{0.873854} & \texttt{0.000321} & \texttt{0.001637} &
\texttt{0.9991,\ 0.0414} \\
\textbf{\texttt{0.275}} & \textbf{16} & \textbf{\texttt{0.49998}} &
\textbf{\texttt{0.276774}} & \textbf{\texttt{0.000000}} &
\textbf{\texttt{0.723226}} & \textbf{\texttt{0.000000}} &
\textbf{\texttt{1.0000,\ 0.0000}} \\
\texttt{0.300} & 16 & \texttt{0.49998} & \texttt{0.061260} &
\texttt{0.000000} & \texttt{0.938740} & \texttt{0.000000} &
\texttt{1.0000,\ 0.0000} \\
\end{longtable}
}

Crucially, this is not mere channel takeover: inside the collapsed site
\(A = 0\), background (\texttt{0.2768}) and particle (\texttt{0.7232})
sectors still coexist, while the radial suppression severs all
connections to the 7 far sites
(\texttt{nz:\ 576\ \$\textbackslash{}to\$\ 16}, \(\|P\| \to 1/2\)).
Collapsing at imaginary site \(A = 1\) (\(\epsilon = 0.30\)) yields the
exact same end-state to 6 digits
(\texttt{E\_bg(1)\ =\ 0.061260,\ E\_P(1)\ =\ 0.938740,\ E\_bg(oth)\ =\ 0.000000}).

\begin{center}\rule{0.5\linewidth}{0.5pt}\end{center}

\subsection{5. Phase-Rigidity Audit vs.~Spin-2 Clock
Dragging}\label{phase-rigidity-audit-vs.-spin-2-clock-dragging}

\begin{enumerate}
\def\labelenumi{\arabic{enumi}.}
\tightlist
\item
  \textbf{Background Phase Rigidity (Audit Finding)}: Although a
  time-averaged site-vector angle ratio initially suggested
  \(\delta\phi/\delta g \approx 0.71\), step-by-step dissection proves
  that \(\delta\phi/\delta g\) fluctuates from \texttt{0.020} to
  \texttt{1.311} with the reference step angle \(\phi_1(\tau)\), and
  direct complex-phase tracking confirms \(|\Delta\theta| < 10^{-9}\)
  rad across all background edges for \(\tau = 1 \dots 100\). Thus, due
  to the rank-1 single-clock rigidity of \(H = O_8 \otimes M_8\),
  cumulative background gravitational frequency shift does not occur in
  weak coupling.
\item
  \textbf{Genuine Particle \& Spin-2 Slowdown}: In contrast, the
  particle sector slows down substantially
  (\(\phi_P(0)/\phi_0 = 0.559\)--\(0.860\)), and the spin-2 candidate
  (\(O_8 \otimes \sigma_x[b=\{2,3\}]\)) locks universally to
  \(d_1 = 0.7071 = 1/\sqrt{2}\) and \(\phi_{b23}/\phi_0 = 0.65778\)
  across \(\epsilon = 0.05, 0.10, 0.20\). At \(\epsilon = 0.20\),
  \(\sigma_x\) drags the entire background clock down to
  \texttt{0.65788} (Fig. 2(b)), and joint P4(\(0.10\)) +
  \(\sigma_x\)(\(0.10\)) injection slows both background
  (\texttt{0.9345\ \$\textbackslash{}to\$\ 0.8664}) and particle
  (\texttt{d1:\ 1.1865\ \$\textbackslash{}to\$\ 0.7602}).
\end{enumerate}

\begin{figure}[H]
\centering
\includegraphics[width=0.96\linewidth,height=\textheight,keepaspectratio,alt={(a) Sharp first-order collapse at \textbackslash epsilon\_c \textbackslash in (0.270, 0.275) (nz: 576 \$\textbackslash to\$ 16, \textbackslash\textbar P\textbackslash\textbar: 0.047 \textbackslash to 1/2, E\_\{\textbackslash mathrm\{bg\}\}(\textbackslash mathrm\{far\}) \textbackslash to 0). (b) Universal spin-2 (\textbackslash sigma\_x) lock (0.6578, d\_1 = 1/\textbackslash sqrt\{2\}) and background clock dragging.}]{figs_en/fig_r5_collapse_and_spin2_en.png}
\caption{(a) Sharp first-order collapse at
\(\epsilon_c \in (0.270, 0.275)\)
(\protect\texttt{nz:\ 576\ \$\textbackslash{}to\$\ 16},
\(\|P\|: 0.047 \to 1/2\), \(E_{\mathrm{bg}}(\mathrm{far}) \to 0\)). (b)
Universal spin-2 (\(\sigma_x\)) lock (\(0.6578, d_1 = 1/\sqrt{2}\)) and
background clock dragging.}
\end{figure}

\begin{center}\rule{0.5\linewidth}{0.5pt}\end{center}

\subsection{6. Discussion, Explicit Limitations, and
Outlook}\label{discussion-explicit-limitations-and-outlook}

This fifth communication establishes four robust geometric/dynamical
phenomena on \(H = O_8 \otimes M_8\)---7-digit isotropic 4-class metric
distortion, discrete curvature gradients, critical far-space
disconnection at \(\epsilon_c \in (0.270, 0.275)\), and spin-2 clock
dragging---while transparently documenting two boundaries: (a)
weak-coupling background phase rigidity (\(|\Delta\theta| < 10^{-9}\)
rad), and (b) the 1-step spatial diameter of the 8-site \(O_8\) cell,
which requires larger spatial lattices or entangled higher layers to
test \(1/r\) potentials, light-cone propagation, and Lorentz factors
\(\sqrt{1-v^2/c^2}\). Rapid Communications 1--5 now form the complete
modular foundation for the unified full paper.

\begin{center}\rule{0.5\linewidth}{0.5pt}\end{center}

\subsection{7. Reproducibility and Released
Files}\label{reproducibility-and-released-files}

\begin{itemize}
\tightlist
\item
  Main sweep: \texttt{tri\_loop\_n64\_gravity\_time\_lorentz\_v0.py}

  \begin{itemize}
  \tightlist
  \item
    Log: \texttt{results/tri\_loop\_n64\_gravity\_time\_lorentz\_v0.txt}
  \end{itemize}
\item
  Falsification audit:
  \texttt{tri\_loop\_n64\_gravity\_rigorous\_check\_v0.py}

  \begin{itemize}
  \tightlist
  \item
    Log:
    \texttt{results/tri\_loop\_n64\_gravity\_rigorous\_check\_v0.txt}
  \end{itemize}
\item
  Edge dissection: \texttt{tri\_loop\_n64\_dphi\_dg\_anatomy\_v0.py}
\item
  Figure generator: \texttt{plot\_fifth\_report\_figures\_v0.py}
\item
  Figures (Fig. 1):
  \texttt{results/fig\_r5\_metric\_and\_curvature\_\{ja,en\}.png}
\item
  Figures (Fig. 2):
  \texttt{results/fig\_r5\_collapse\_and\_spin2\_\{ja,en\}.png}
\item
  TeX builder: \texttt{build\_tex\_rapid5.py}
\end{itemize}

\subsection{References}\label{references}

\begin{enumerate}
\def\labelenumi{\arabic{enumi}.}
\tightlist
\item
  N. Kihara, \emph{Rapid Communication 1}, Zenodo (2026), DOI:
  10.5281/zenodo.24616931.
\item
  N. Kihara, \emph{Second Rapid Communication}, Zenodo (2026), DOI:
  10.5281/zenodo.24618400.
\item
  N. Kihara, \emph{Third Rapid Communication}, Zenodo (2026), DOI:
  10.5281/zenodo.24619357.
\item
  N. Kihara, \emph{Fourth Rapid Communication}, Zenodo (2026), DOI:
  10.5281/zenodo.24620262.
\end{enumerate}

\end{document}
